\documentclass[12pt,letterpaper]{article}

\usepackage[utf8]{inputenc}
\usepackage[spanish]{babel}
\usepackage{graphicx}
\usepackage{geometry}
\usepackage{amsmath}
\usepackage{float}
\usepackage{caption}
\usepackage{subcaption}
\usepackage{setspace}
\usepackage{csquotes}
\usepackage{booktabs}
\usepackage{titling}
\usepackage{siunitx}
\usepackage[backend=biber,style=ieee]{biblatex}
\addbibresource{references.bib}

\geometry{left=2.5cm, right=2.5cm, top=2.5cm, bottom=2.5cm}

\graphicspath{{figuras/}}

% --- Figura con placeholder automático: compila con o sin la imagen ----------
% Uso: \figph{archivo.png}{Pie de figura}{clave_label}
% Si figuras/archivo.png no existe, dibuja un recuadro PLACEHOLDER.
\usepackage{xcolor}
\newcommand{\figph}[3]{%
  \begin{figure}[H]
    \centering
    \IfFileExists{figuras/#1}{%
      \includegraphics[width=0.9\textwidth]{#1}%
    }{%
      \fcolorbox{gray}{gray!8}{\begin{minipage}[c][0.22\textheight][c]{0.88\textwidth}%
        \centering\footnotesize\color{gray!55!black}%
        [\,PLACEHOLDER\,]\\[4pt] \texttt{\detokenize{figuras/#1}}%
      \end{minipage}}%
    }
    \caption{#2}
    \label{fig:#3}
  \end{figure}%
}

% =============================================================================
% Reporte de Avance — Cuarto Semestre
% Continuidad del reporte de Tercer Semestre (instrumentación PA propietaria).
% Estado: BORRADOR. Las marcas % TODO indican datos/figuras pendientes del autor.
% =============================================================================

\title{Análisis de Tejidos mediante Tomografía Fotoacústica e IA\\
Validación de la Cadena de Adquisición con ADC Físico y Diagnóstico del
Front-End de Detección Ultrasónica}
\author{Eduardo Santos Mena \\ Asesores: Dr. Samuel Pérez Huerta, Dr. Augusto David Ariza Flores \\
\small Unidad Académica de Ciencia y Tecnología de la Luz y la Materia, UAZ}
\date{Reporte de Avance - Cuarto Semestre 2026} % TODO: confirmar periodo/fechas

\begin{document}

\maketitle

\begin{abstract}
\noindent
Este reporte da continuidad a la línea de ingeniería de instrumentación
fotoacústica (PAT) iniciada en el periodo anterior, en el que el núcleo digital se
validó con señales sintéticas (\textit{Virtual ADC}). En este cuarto periodo se
completó la integración física del conversor analógico-digital (AD9226) con la
FPGA y se certificó, ahora con señal real, que la cadena de adquisición de datos
(DAQ) captura, almacena y transmite sin pérdida de información, alcanzando una
resolución de 12 bits a una tasa de muestreo de 54~MSPS. Asimismo, se consolidó la
operación del Driver de excitación V2.0 con retroalimentación óptica. Al llevar el
sistema completo al banco para la primera adquisición fotoacústica, se obtuvo un
\textbf{resultado negativo rigurosamente controlado}: la señal observada,
sincronizada con el disparo láser, se demostró mediante tres controles
independientes (desconexión del sensor, ausencia de retardo de tiempo de vuelo e
independencia respecto a la distancia) como \textbf{acoplamiento electromagnético
(pickup)}, no como detección acústica. El análisis localiza la barrera en el primer
eslabón de la cadena de detección —un transductor piezoeléctrico de alta impedancia
conectado en crudo al digitalizador, sin pre-amplificación— y define el front-end
analógico de detección (pre-amplificador de bajo ruido, acoplamiento de impedancias,
apantallamiento y sincronía por disparo óptico) como el paso decisivo. Se propone un
protocolo de \textit{bring-up} por capas de la detección, un plan de contingencia y
una re-planificación del cronograma.

En conjunto, este periodo cierra la validación del subsistema digital y aísla con
precisión el cuello de botella real del sistema, trasladando el reto desde el
dominio del cómputo (resuelto) hacia el de la detección de señales débiles en
presencia de una excitación de alto voltaje.
\end{abstract}

\newpage
\section{Introducción}

La tomografía fotoacústica (PAT) es una modalidad de imagen híbrida que combina la
especificidad molecular de la excitación óptica con la resolución de la detección
ultrasónica, sin radiación ionizante~\cite{lin_emerging_2022}. Este trabajo se
enmarca en un proyecto de doctorado (LUMAT--UAZ) cuyo objetivo general es
\textbf{desarrollar y evaluar un sistema integrado de PAT y \textit{deep learning}
para la detección temprana de la artritis reumatoide (AR)}, combinando
instrumentación propia de bajo costo y algoritmos de reconstrucción y segmentación
basados en redes neuronales convolucionales~\cite{hauptmann_deep_2020}
~\cite{grohl_deep_2021}. La AR es un caso de estudio idóneo: la PAT puede revelar
cambios vasculares precoces (angiogénesis sinovial) antes de que sean visibles en
las modalidades convencionales. El protocolo plantea como metas del sistema una
resolución espacial de 150--200~$\mu$m y una profundidad efectiva de 1--2~cm en
tejido articular.

La estrategia del proyecto se divide en dos etapas: (a) diseño y validación del
sistema de adquisición PAT, y (b) desarrollo de los algoritmos de análisis. El
esfuerzo de instrumentación ha progresado de forma incremental: en el \textbf{segundo
semestre} se diseñó y validó la fuente de excitación óptica —un driver de corriente
en modo avalancha que dispara láseres LED de \SI{445}{\nano\meter} con pulsos
inferiores a \SI{30}{\nano\second}, a costo accesible~\cite{wu_advanced_2021}
~\cite{li_photoacoustic_2020}—; en el \textbf{tercer semestre} se evolucionó dicha
fuente (Driver V2.0) y se desarrolló el núcleo digital del sistema de adquisición,
certificado con señales sintéticas internas (\textit{Virtual ADC}).

El presente periodo aborda la transición desde la validación sintética hacia la
operación con \textbf{hardware físico completo}: la integración del conversor
analógico-digital con la lógica de la FPGA y el primer intento de adquisición de
una señal fotoacústica real. Como se documenta en este reporte, dicho intento
reveló que, una vez resuelto el subsistema digital, la dificultad dominante reside
en la \textbf{interfaz analógica de detección}: la conversión de una onda acústica
débil, captada por un transductor piezoeléctrico de alta impedancia, en una señal
eléctrica observable, en un entorno donde la excitación de alto voltaje constituye
un agresor electromagnético sincronizado con la medición.

\section{Objetivos del Periodo}

El objetivo general fue validar la cadena de adquisición completa con hardware
físico y realizar la primera caracterización de la cadena de detección
fotoacústica.

Objetivos específicos:
\begin{enumerate}
    \item \textbf{Integración física ADC--FPGA:} montar el módulo AD9226 y
    certificar, con señal real, la integridad, linealidad y respuesta dinámica de
    la cadena de adquisición.
    \item \textbf{Operación del Driver V2.0:} verificar la estabilidad de la fuente
    de excitación y su subsistema de retroalimentación óptica.
    \item \textbf{Primera adquisición de detección:} intentar la captura de una
    señal acústica/fotoacústica y caracterizar el desempeño de la cadena de
    detección.
\end{enumerate}

\section{Metodología y Desarrollo}

\subsection{Integración física del ADC y validación por etapas}
La integración del digitalizador se realizó siguiendo una metodología de
\textit{bring-up} por etapas, en la que cada capa valida una sola variable y no se
avanza hasta que la anterior se comporta exactamente como se espera. Esta
estrategia, heredada de la depuración del enlace digital, permite atribuir
cualquier fallo a lo que la etapa en curso agregó.

El módulo AD9226 se montó sobre la plataforma de integración del sistema: una placa
de acrílico que aloja, atornilladas, las PCBs de todas las etapas
(Figura~\ref{fig:acrilico}).

\figph{placa-acrilico.png}{Plataforma de integración: placa de acrílico con las PCBs
de todas las etapas del sistema atornilladas.}{acrilico}

\subsection{Caracterización del DAQ en rango}
Con el ADC físico operando, se repitieron los protocolos de validación para
confirmar el comportamiento de la cadena completa: integridad de transmisión,
linealidad de la conversión y reconstrucción de señales dinámicas, ahora sobre
datos reales y no sintéticos.

\subsection{Montaje de detección y protocolo de adquisición}
El transductor piezoeléctrico (Yushi, \SI{2.5}{\mega\hertz}) se conectó mediante un
conector BNC y un cable BNC--SMA directamente a la entrada del módulo AD9226 y,
en paralelo, al osciloscopio para observación directa. La excitación se realizó con
el Driver V2.0. Para discriminar el origen de la señal observada se diseñaron
controles de descarte (desconexión del sensor, variación de la distancia y análisis
del retardo temporal), descritos en la Sección~\ref{sec:resultados}.

\figph{montaje.png}{Montaje de banco empleado en el intento de adquisición de
detección.}{montaje}

\subsection{Driver V2.0 en operación}
El Driver V2.0 operó conforme a lo establecido en el periodo anterior: control
embebido (RP2040), gestión térmica y retroalimentación óptica mediante fotodiodo y
amplificador de transimpedancia para la estimación de energía por pulso.

\section{Resultados}
\label{sec:resultados}

\subsection{Cadena de adquisición con ADC físico (positivo)}
La integración física del AD9226 con la FPGA quedó validada: el sistema captura,
almacena en memoria interna y transmite las ráfagas sin pérdida de datos,
operando a \SI{54}{\mega\hertz} (54~MSPS) con una resolución
efectiva de 12 bits (LSB $\approx \SI{2.44}{\milli\volt}$). La resolución temporal
resultante (aproximadamente 21.6 muestras por ciclo a \SI{2.5}{\mega\hertz}) es
suficiente para resolver el \textit{ringing} característico del transductor.

\figph{bringup-etapas.png}{Validación por etapas del FPGA con el ADC físico:
resultados de los experimentos de integridad, linealidad y respuesta dinámica
(metodología \textit{bring-up}, carpeta \texttt{bringup/}).}{bringup}

\subsection{Driver de excitación V2.0 (positivo)}
El Driver V2.0 entregó pulsos de corriente estables (Figura~\ref{fig:pulsos}),
manteniendo la operación de la fuente de excitación durante las pruebas.

\figph{pulsos-corriente.png}{Pulsos de corriente del Driver V2.0.}{pulsos}

\subsection{Primera adquisición de detección: resultado negativo controlado}
Al excitar el sistema, el canal del transductor (CH2) registró una onda de
aproximadamente \SI{150}{\milli\volt} pico a pico, coincidente con el pulso láser
(CH1) y con apariencia de \textit{ringing} en torno a \SI{2.6}{\mega\hertz}. Tres
controles independientes demuestran que esta señal \textbf{no es de origen
acústico}:

\begin{enumerate}
    \item \textbf{Desconexión del sensor:} la misma onda persiste con el transductor
    físicamente desconectado, por lo que no proviene del elemento piezoeléctrico.
    \item \textbf{Ausencia de retardo (t = 0):} la señal coincide exactamente con el
    disparo láser, sin el retardo de tiempo de vuelo que tendría una onda acústica;
    es, por tanto, de naturaleza electromagnética.
    \item \textbf{Independencia de la distancia:} variar la distancia entre la fuente
    y el sensor no modifica la señal, confirmando la ausencia de dependencia con el
    tiempo de vuelo.
\end{enumerate}

Se concluye que la señal observada corresponde a \textbf{acoplamiento
electromagnético (pickup)} sincronizado con la descarga de excitación, y no a una
detección acústica. La frecuencia de \textit{ringing} observada queda determinada
por la red de entrada (cable y impedancia de entrada del digitalizador) y no por la
resonancia del transductor, como evidencia su persistencia con el sensor
desconectado. Una prueba con fuente mecánica conocida (rotura de mina,
Hsu--Nielsen, sin láser) tampoco produjo señal observable, resultado consistente
con la ausencia de pre-amplificación y con el hecho de que dicha fuente concentra su
energía por debajo de \SI{1}{\mega\hertz}, mal acoplada a un transductor de
\SI{2.5}{\mega\hertz}.

\figph{ringing-falso.png}{Control de descarte. CH1: pulso láser; CH2: onda de
$\sim$150~mV pp coincidente con el disparo y con apariencia de \textit{ringing}
($\sim$2.6~MHz). La señal persiste con el sensor desconectado y llega en $t=0$
(sin retardo de tiempo de vuelo), evidenciando acoplamiento electromagnético
(\textit{pickup}) y no detección acústica.}{pickup}

\section{Discusión}

\subsection{El cuello de botella está en el front-end de detección}
Los resultados separan con claridad los subsistemas validados (ADC, FPGA, enlace,
Driver V2.0) del subsistema bloqueante. Como ni una fuente mecánica conocida ni el
láser produjeron señal acústica observable, la incógnita no reside en el láser ni en
el dominio digital, sino en la cadena
sensor~$\rightarrow$~acoplamiento~$\rightarrow$~(pre)amplificación~$\rightarrow$~observabilidad.
La causa principal identificada es la \textbf{ausencia de un pre-amplificador}: el
transductor, fuente de alta impedancia, se conecta en crudo al digitalizador y al
osciloscopio, de modo que una señal fotoacústica esperada en el orden de los
microvoltios a milivoltios queda por debajo tanto del piso de observación como del
\textit{pickup} de \SI{150}{\milli\volt}.

\subsection{Protocolo de \textit{bring-up} de la detección}
Se propone validar la cadena de detección por capas, en analogía con el
\textit{bring-up} digital:
\begin{itemize}
    \item \textbf{D0 — Salud eléctrica del piezo:} medición de capacitancia (LCR) y
    continuidad para confirmar que el elemento está intacto.
    \item \textbf{D1 — Respuesta mecánica directa} (con el láser apagado): un
    estímulo mecánico en la cara del transductor debe producir un transitorio
    observable en el osciloscopio, sin posibilidad de confusión con \textit{pickup}.
    \item \textbf{D2 — Fuente acústica conocida y acoplada}, preferentemente
    \textit{in-band}.
    \item \textbf{D3 — Incorporación del pre-amplificador} de bajo ruido junto al
    sensor.
    \item \textbf{D4 — Reintroducción del láser}, buscando la señal en la ventana
    retardada por tiempo de vuelo y verificando su desplazamiento con la distancia.
\end{itemize}

\subsection{Especificación del pre-amplificador}
El pre-amplificador debe amplificar señales débiles \textbf{en el sensor}, sin
cargar el transductor de alta impedancia y sin añadir ruido apreciable, sobre la
banda del transductor. Las especificaciones objetivo son: impedancia de entrada
$\geq \SI{1}{\mega\ohm}$; ganancia de 40--60~dB; ancho de banda plano de
0.1--5~MHz; ruido referido a la entrada $\leq$ 3--5~nV/$\sqrt{\text{Hz}}$; y
apantallamiento integral. Una salida de baja impedancia, además, reduce de forma
sustancial la susceptibilidad al \textit{pickup} en el cableado respecto a la línea
cruda de alta impedancia.

% TODO: tabla de trade-off de candidatos (p. ej. AD8331/AD8332 de ultrasonido) y
% selección final.

\subsection{El acoplamiento electromagnético como artefacto a controlar}
El \textit{pickup} de la descarga de alto voltaje no es la barrera principal, pero
debe controlarse para no confundirlo con señal y para evitar que sature la entrada
cuando exista detección real. Las medidas previstas son: sincronía por disparo
óptico (empleando el fotodiodo/TIA del Driver V2.0 como referencia temporal limpia),
apantallamiento tipo Faraday, y la separación temporal de la señal (ventana
retardada o pre-trigger), dado que la señal acústica real llega con un retardo de
tiempo de vuelo respecto al \textit{pickup} en $t = 0$.

\section{Limitaciones}
\begin{itemize}
    \item La validación de la detección está pendiente; aún no se ha demostrado una
    señal acústica real por encima del piso de observación.
    \item La ventana de captura actual (5~$\mu$s, $\approx \SI{7.7}{\milli\meter}$ de
    profundidad) es insuficiente para la cobertura de tejido objetivo; su ampliación
    requiere la implementación de \textit{buffers} (FIFO) en la FPGA.
    \item No se dispone de \textit{pulser} ni de un segundo transductor para una
    validación acústica \textit{in-band} inmediata (véase Plan~B/C).
\end{itemize}

\section{Conclusiones}
\begin{enumerate}
    \item \textbf{Subsistema digital cerrado:} la cadena de adquisición con ADC
    físico quedó validada (12 bits, 54~MSPS, sin pérdida de datos), superando lo
    proyectado para el periodo.
    \item \textbf{Cuello de botella aislado:} mediante controles independientes se
    demostró que la lectura observada es acoplamiento electromagnético y no
    detección acústica, localizando la barrera en el front-end analógico de
    detección.
    \item \textbf{Camino accionable:} se define un protocolo de \textit{bring-up} de
    la detección y la especificación de un pre-amplificador como paso decisivo.
\end{enumerate}

\section{Plan de Contingencia (Plan B / C)}
Dado que no se dispone de \textit{pulser} ni de un segundo transductor, la validación
acústica \textit{in-band} se prioriza por viabilidad:
\begin{itemize}
    \item \textbf{B-1 (equipo actual):} láser sobre un absorbente fuerte (grafito,
    tinta china) acoplado por gel a distancia conocida, buscando la señal en la
    ventana retardada por tiempo de vuelo y verificando su desplazamiento con la
    distancia.
    \item \textbf{B-2:} pulso-eco con generador de funciones (\textit{burst} de
    \SI{2.5}{\mega\hertz}) e inyección eléctrica calibrada para caracterizar la
    cadena receptora.
    \item \textbf{C:} adquisición/fabricación de un \textit{pulser} de ultrasonido y/o
    un segundo transductor para validación por pulso-eco / transmisión en tanque de
    agua.
    \item \textbf{Contingencia de alcance:} de persistir la dificultad, reencuadrar
    el entregable como plataforma de instrumentación PA de bajo costo con cadena de
    recepción caracterizada, difiriendo la imagen en fantoma.
\end{itemize}

\section{Re-planificación y Evaluación del Retraso}

El retraso está \textbf{localizado en un único hito} (la primera adquisición
acústica) y no afecta el camino de datos, ya cerrado. Su magnitud está acotada al
ciclo de diseño, fabricación y prueba del front-end analógico; las etapas
posteriores (procesamiento de señal e IA) son predominantemente de software sobre
datos ya capturables, por lo que pueden comprimirse una vez establecida la
detección.

\begin{table}[H]
    \centering
    \begin{tabular}{|p{0.16\textwidth}|p{0.40\textwidth}|p{0.36\textwidth}|}
    \hline
    \textbf{Bloque} & \textbf{Actividad principal} & \textbf{Entregable} \\ \hline
    B1 & Diseño del front-end de detección: pre-amplificador, acoplamiento de
    impedancias, apantallamiento y sincronía óptica & Esquemático + simulación +
    presupuesto de ruido \\ \hline
    B1$'$ & Validación de la cadena de recepción (Plan~B-1/B-2) & Señal acústica
    real capturada (de-riesgo) \\ \hline
    B2 & Fabricación y prueba del front-end & Cadena de recepción caracterizada
    (ganancia/BW/ruido) \\ \hline
    B3 & Firmware: ventana de captura ampliada (FIFO) y ventana retardada/pre-trigger
    & A-scan de profundidad útil \\ \hline
    B4 & Primera adquisición PA en fantoma (gelatina + grafito) & A-scan PA con
    tiempo de vuelo coherente \\ \hline
    B5 & Procesamiento de señal (retroproyección) & Reconstrucción básica \\ \hline
    B6 & IA preliminar (denoising 1D) y escritura & Artículo / reporte \\ \hline
    \end{tabular}
    \caption{Cronograma revisado.} % TODO: asignar meses/fechas reales del periodo.
    \label{tab:cronograma_rev}
\end{table}

\section{Reflexiones}
La metodología de validación por capas permitió afirmar con certeza que el
subsistema digital no es el problema y señalar el analógico: ese aislamiento es, en
sí mismo, el avance del periodo. El proyecto transita de la pregunta
\enquote{¿podemos digitalizar?} (resuelta) a \enquote{¿podemos detectar señales
débiles junto a una excitación de alto voltaje?}, que es la pregunta correcta y más
difícil. La lección —que en sistemas fotoacústicos de bajo costo el front-end
analógico y la integridad electromagnética dominan el éxito tanto como el cómputo—
constituye una contribución relevante para el artículo de instrumentación previsto.

\printbibliography

\end{document}
